\section*{Глоссарий: DEIC и ключевые понятия}
\label{sec:glossary}
\addcontentsline{toc}{section}{Глоссарий}

Статья писалась на пересечении психометрии, клинической психологии и контемплативных традиций. Чтобы читатель не застревал на технических терминах, здесь собраны краткие определения. Первые упоминания каждого термина в тексте дублируются с русским эквивалентом; в таблицах и формулах остаётся сокращённое обозначение.

\subsection*{Центральные конструкты DEIC}

\begin{description}
    \item[DEIC] \textit{Dynamic Empathy and Inner Conflict} --- «динамическая эмпатия и внутренний конфликт»; название измерительного каркаса и одноимённого опросника из 131 пункта.
    \item[ID (Internal Dissonance)] \textbf{внутренний диссонанс} --- состояние, когда эмпатический отклик уже пошёл, но дальше блокируется из-за невыносимости аффективного материала \emph{как такового}, --- чужого или собственного, эти два направления в ID симметричны и есть одна и та же контейнирующая функция. Переживается как «хочу чувствовать, но не могу»: не только к другому, но и к себе. Пассивная защита, низкое «разрешение» чувствования в обе стороны одновременно.
    \item[CA (Childhood Alienation)] \textbf{детское отчуждение} --- ранняя семейно-эмоциональная история: отсутствие отклика значимых взрослых, невозможность быть собой, переживание небезопасности ближайшего круга. Не сумма травмирующих событий, а климат.
    \item[P (Psychopathy)] \textbf{психопатический паттерн} --- архитектура, в которой эмпатия \emph{не заблокирована}, а \emph{условна}: я чувствую того, кого выбираю чувствовать, и могу «переключить» этот канал по желанию. В DEIC --- морально нейтральная координата, не диагноз.
    \item[M (Masking)] \textbf{маскирование} --- социальное предъявление, не совпадающее с внутренним состоянием. В текущей версии инструмента шкала слабая ($\omega = 0.55$); в волне 2 будет переработана.
    \item[E\textsubscript{AF} (Affective Empathy)] \textbf{аффективная эмпатия} --- способность к непосредственному эмоциональному резонансу: чужая боль отзывается в теле как своя.
    \item[E\textsubscript{cog} (Cognitive Empathy)] \textbf{когнитивная эмпатия} --- способность понять внутреннее состояние другого как модель, как ментальную карту, без обязательного телесного резонанса. Это та эмпатия, которую можно обучить и которую используют следователь, психопат и хороший терапевт --- каждый по-своему.
    \item[A (Attention / Witnessing)] \textbf{внимание, присутствие-свидетель} --- метакогнитивная способность наблюдать собственные состояния, не сливаясь с ними. Ближе всего к тому, что контемплативные традиции называли satipa\d{t}\d{t}h\=ana (буддизм), s\=ak\d{s}\=i (адвайта), rigpa (дзогчен), self-remembering (Гурджиев). В DEIC --- не содержание, а условие наблюдения.
\end{description}

\subsection*{Психометрические методы}

\begin{description}
    \item[Фактор (латентный)] --- предполагаемая внутренняя характеристика, которую нельзя измерить напрямую, но можно реконструировать по совокупности ответов на пункты опросника. Например, «добросовестность» --- это не один вопрос, а общее, что объединяет ответы на десяток пунктов.
    \item[EFA] \textit{exploratory factor analysis} --- \textbf{разведочный факторный анализ}: «сколько скрытых характеристик достаточно, чтобы объяснить структуру ответов?» Применяется, когда теория ещё не фиксирует точное число факторов.
    \item[CFA] \textit{confirmatory factor analysis} --- \textbf{подтверждающий факторный анализ}: «действительно ли пункты группируются так, как теория предсказывает?» Проверка, а не разведка.
    \item[SEM] \textit{structural equation modeling} --- \textbf{структурное моделирование}: метод, в котором несколько регрессий оцениваются одновременно, с латентными переменными и путями между ними. Позволяет спросить, например: «предиктор A влияет на исход C напрямую или только через медиатор B?» --- и получить ответ в единой системе уравнений.
    \item[Рефлексивный латент] --- классическая модель: скрытый фактор \emph{порождает} наблюдаемые пункты; изменение уровня фактора отражается во всех пунктах одновременно. Пункты взаимозаменяемы как индикаторы.
    \item[Формативный композит] --- противоположная конструкция: набор разнородных пунктов \emph{складывается} в итоговую величину, не отражающую единое скрытое свойство. Здесь удаление одного пункта меняет смысл композита. A в DEIC организован именно так.
    \item[Факторный скор] --- расчётное значение латентного фактора для каждого участника, полученное из CFA; используется как переменная в последующих анализах.
    \item[$\omega$ (омега МакДональда)] --- современная мера внутренней согласованности шкалы (надёжности), предпочтительная альтернатива Кронбаха. Значения выше $0.80$ --- приемлемо, выше $0.90$ --- высокая надёжность.
    \item[Parcels (парселы)] --- агрегированные мини-композиты из 3--5 пунктов каждый, используемые вместо сырых айтемов в CFA. Уменьшают число параметров модели и сглаживают шум айтем-уровня.
    \item[Двухуровневое scoring (DEIC wave-2)] --- предлагаемая архитектура опросника, в которой каждый айтем имеет \emph{два} технических поля: \texttt{factor\_weights} (вектор весов по нескольким факторам, используемый в CFA/SEM и сохраняющий теоретическое многомерное содержание пункта) и \texttt{primary\_scale} (единственный фактор, в который айтем входит при композитном sum-scoring). Разделение устраняет алгебраическое раздувание корреляций между композитами, характерное для наивного multi-scale scoring, при сохранении multi-scale весов как \emph{теоретической} а-priori-декларации. Отличается от «один айтем $=$ один фактор» тем, что cross-loadings остаются в модели, а не вычёркиваются.
\end{description}

\subsection*{Причинно-следственный анализ путей}

\begin{description}
    \item[Медиация (опосредование)] --- ситуация, когда A влияет на C не напрямую, а \emph{через} промежуточную переменную B. Пример: детское отчуждение влияет на аффективную эмпатию не непосредственно, а через формирование внутреннего диссонанса.
    \item[a-путь] --- звено A $\to$ B (предиктор воздействует на медиатор).
    \item[b-путь] --- звено B $\to$ C (медиатор воздействует на исход).
    \item[Прямой эффект (c$'$)] --- часть влияния A на C, остающаяся после того, как медиатор учтён.
    \item[Непрямой эффект] --- часть влияния A на C, проходящая через медиатор (произведение a $\cdot$ b).
    \item[Модерация] --- ситуация «при каких условиях»: эффект A на B не постоянный, а зависит от уровня третьей переменной M. В DEIC A модерирует и a-путь (CA $\to$ ID), и b-путь (ID $\to$ E\textsubscript{AF}).
    \item[Взаимодействие (A $\times$ B)] --- формальное выражение модерации в регрессии: произведение двух переменных, коэффициент при котором говорит, насколько их общий эффект нелинеен.
    \item[IMM] \textit{index of moderated-moderated mediation} --- \textbf{индекс двойной модерации}: произведение двух коэффициентов взаимодействия ($a_{\text{mod}} \times b_{\text{mod}}$). Если IMM значимо отличается от нуля, модератор действует на \emph{всю цепочку} A$\to$B$\to$C одновременно, а не на отдельное её звено. В DEIC отрицательный IMM означает: высокое внимание разрывает каскад «травма $\to$ диссонанс $\to$ блок эмпатии» на обоих концах сразу.
    \item[Bootstrap] --- процедура многократного пересемплирования данных с возвращением (обычно 2000--5000 повторений) для оценки доверительных интервалов сложных коэффициентов, когда формульное решение невозможно.
    \item[95\,\% CI (доверительный интервал)] --- диапазон, в который попадает истинное значение коэффициента с вероятностью 95\,\%. Если CI не включает ноль, эффект значим на уровне $p < 0.05$.
\end{description}

\subsection*{Индексы пригодности модели}

\begin{description}
    \item[$\chi^2$ (хи-квадрат)] --- статистика несоответствия модели данным. Сама по себе у больших выборок почти всегда значима, поэтому оценивается в паре с приведёнными ниже индексами.
    \item[df (степени свободы)] --- число независимых ограничений, по которым проверяется модель; чем больше df, тем «экономнее» модель.
    \item[CFI] \textit{comparative fit index} --- \textbf{сравнительный индекс соответствия}: $>0.90$ приемлемо, $>0.95$ хорошо.
    \item[TLI] \textit{Tucker--Lewis index} --- родственен CFI, чуть строже к сложности модели.
    \item[RMSEA] \textit{root mean square error of approximation} --- \textbf{средняя ошибка аппроксимации}: $<0.06$ хорошо, $<0.08$ приемлемо, $>0.10$ плохо.
    \item[Satorra--Bentler scaling] --- поправка хи-квадрата на ненормальность распределения данных; особенно важна для Likert-шкал.
    \item[MLR] \textit{robust maximum likelihood} --- метод оценки со стандартными ошибками, устойчивыми к ненормальности. Reviewer-стандарт для опросниковых данных; требует R/lavaan.
\end{description}

\subsection*{Клинические и теоретические понятия}

\begin{description}
    \item[ACE] \textit{adverse childhood experiences} --- шкала неблагоприятного детского опыта; в DEIC используется как ковариата и как компонент композитов.
    \item[Первичная психопатия (Karpman, 1941)] --- тип с низким внутренним диссонансом и без выраженной детской травмы; «конституциональный» холод, минимальная эмоциональная реактивность.
    \item[Вторичная психопатия] --- тот же психопатический паттерн переключаемой эмпатии, но с высоким ID и высоким CA; «травма-порождённый» путь, принципиально иной терапевтический профиль.
    \item[Integrated survivor] \textbf{интегрированный выживший} --- профиль с высоким CA (тяжёлое детство) и высоким A (развитая работа с вниманием); в нашей выборке имеет \emph{наивысшую} аффективную эмпатию. Эмпирическая опора тезиса «травма не приговор».
    \item[Hyper-instrumentation] \textbf{гиперинструментализация} --- гипотеза (Blair, Baron-Cohen): психопатический паттерн не \emph{понижает}, а \emph{обостряет} когнитивную эмпатию, используя её как инструмент. В DEIC подтверждается положительной частной корреляцией $r(\text{P}, \text{E}_{\text{cog}}) = +0.32$ после контроля общих источников дисперсии.
    \item[Свобода воли (операциональная)] --- в DEIC не бинарное метафизическое свойство и не иллюзия, а \emph{градуальная операциональная переменная}: величина зазора между автоматическим стимул-реакция-связыванием (работой D/P/E-конфигурации) и наблюдением за этим связыванием (A). Нулевая A $\Rightarrow$ поведение функционально детерминировано текущей конфигурацией; высокая A $\Rightarrow$ увеличивается пространство выбора, пропорционально моментальной способности удерживать наблюдающее отношение. Это операционализация того, что традиции описывают как «пробуждение», «освобождение», «раб воли», «самовоспоминание»; не решение метафизической проблемы свободы воли, а её перефокусировка в измеримую психологическую переменную.
\end{description}

\subsection*{Контемплативные ссылки}

\begin{description}
    \item[Satipa\d{t}\d{t}h\=ana] (палийский буддизм) --- «установление внимательности»; систематическая практика наблюдения тела, чувств, состояний ума и феноменов.
    \item[S\=ak\d{s}\=i] (адвайта-веданта) --- «свидетель»; то в опыте, что остаётся неизменным при смене содержаний.
    \item[Rigpa] (тибетский дзогчен) --- знание-присутствие, не-двойственное осознавание, предшествующее разделению на субъект и объект.
    \item[Self-remembering] (Гурджиев, четвёртый путь) --- практика удерживания осознавания себя параллельно с вовлечённостью в действие; ближайший западный эквивалент witnessing.
    \item[Anti-Wilber (дисциплинарное правило)] --- отказ от иерархической «master-map», накладываемой поверх традиций. DEIC позиционируется как операциональный диалект того, что традиции видели, а не как мета-карта над ними.
\end{description}

\subsection*{Нейронаучные отсылки}

\begin{description}
    \item[DMN] \textit{default mode network} --- «сеть пассивного режима работы мозга»; активна при блуждающем уме и самореферентных процессах. Гипотетический нейронный коррелят того, что A в DEIC меряет на феноменологическом уровне.
    \item[Predictive coding (предиктивное кодирование)] --- нейронаучная рамка, в которой мозг работает как машина, непрерывно предсказывающая сенсорный ввод и корректирующая модель по ошибкам прогноза. В статье используется как кандидатная рамка для объяснения различия ID (bottom-up перегруз) и P (top-down выбор модели).
\end{description}

\clearpage
